\section{Background}
Listening to a machine is one of the oldest ways of judging its health. An experienced operator can tell from a change in pitch or rhythm that a bearing is wearing or a valve is sticking, long before a sensor on the machine reports a fault. Microphones are inexpensive, need no contact with the machine and can listen continuously, which makes acoustic condition monitoring attractive for factories with many small machines. The difficulty is on the learning side. Failures are rare, their sounds are diverse, and the failure that matters most is usually one that has never been recorded. A detector therefore cannot be trained as a classifier of known faults; it has to learn what \emph{normal} sounds like and raise an alarm when a recording departs from it. This is the setting of unsupervised anomalous sound detection (ASD), which the DCASE Challenge has run as Task~2 since 2020 \cite{koizumi2020,dohi2023}.

The 2026 edition of the task \cite{dcase2026task2} is described as noise-aware. Each clip is recorded with two channels, and every machine type comes in two domains: a \emph{source} domain with plenty of normal training clips, and a \emph{target} domain, recorded under different conditions or background noise, for which only ten normal training clips are available. A useful detector must work in both domains even though it has seen almost nothing of the second one. This domain shift, rather than the modelling of normal sound itself, turned out to be the central difficulty in our results.

\section{What this project does}
We implemented a complete detector for all seven machine types of the development dataset, following the experimental plan given in the AI Laboratory problem statement:
\begin{enumerate}
  \item \textbf{Experiment~1} trains an Autoencoder on normal recordings and scores a test recording by its reconstruction error.
  \item \textbf{Experiment~2} replaces the reconstruction error by a Mahalanobis distance from the distribution of normal features, computed both in the Autoencoder's latent space and in its residual space, and compares the two scoring approaches.
  \item \textbf{Experiment~3} trains the same Autoencoder twice per machine, once on source-domain recordings only (Model~A) and once on source plus target recordings (Model~B), to measure what the ten target clips contribute.
\end{enumerate}
The work was split between the two of us. \studentOne{} (Part~1) audited the data, built the preprocessing and the Autoencoder, trained the fourteen models in Google Colab and prepared the statistics and handover package. \studentTwo{} (Part~2) scored the test recordings, calibrated thresholds, computed and compared the metrics and built the Streamlit application. Both parts share one library module, so training, evaluation and the application run byte-identical preprocessing and model code.

\section{How to read this report}
Chapters~\ref{ch:problem} to \ref{ch:env} describe the problem, the data as we actually found it, and the method and its implementation. Chapter~\ref{ch:results} presents the results, recomputed during a final audit from the per-recording score files. Chapters~\ref{ch:app} to \ref{ch:conclusion} describe the application, discuss what the results mean, state the limitations plainly and suggest what we would do next. The appendix lists folder structures, configuration values and the exact code interface used by Part~2.
